% Generated draft: requires mathematical and source review.
\documentclass[11pt]{article}
\usepackage[margin=25mm]{geometry}
\usepackage{fontspec}
\setmainfont{Latin Modern Roman}
\usepackage{amsmath,amssymb,mathtools,unicode-math}
\setmathfont{Latin Modern Math}
\usepackage{enumitem,needspace,xurl,hyperref}
\setlength{\parindent}{0pt}
\setlength{\parskip}{5pt}
\setlength{\emergencystretch}{3em}
\newcommand{\N}{\mathbb{N}}
\newcommand{\Z}{\mathbb{Z}}
\newcommand{\Q}{\mathbb{Q}}
\newcommand{\R}{\mathbb{R}}
\newcommand{\C}{\mathbb{C}}
\newcommand{\F}{\mathbb{F}}
\newcommand{\A}{\mathbb{A}}
\newcommand{\PP}{\mathbb{P}}
\newcommand{\E}{\mathbb{E}}
\newcommand{\eps}{\varepsilon}
\newcommand{\dd}{\,\mathrm{d}}
\begin{document}
\section*{UnsolvedMath ID: 11000045}
\textbf{Source problem code:} catalogue entry
\subsection{Definitions and mathematical statement}
Let $\mathcal M_g=\operatorname{Mod}_g\backslash\mathcal T_g$ be the moduli space of closed genus-$g$ Riemann surfaces, $g\ge2$, with the quotient Teichm\"uller metric $d_{\mathrm{Teich}}$. The immediately preceding source record, local ID 11000044, defines its tangent cone at infinity by the pointed Gromov--Hausdorff limit
\[
 \operatorname{Cone}(\mathcal M_g)=
 \lim_{n\to\infty}(\mathcal M_g,n^{-1}d_{\mathrm{Teich}},x_0),
\]
for a fixed basepoint $x_0$, and reports that this limit exists in this setting. This is a large-scale metric cone construction, not the Deligne--Mumford compactification or its boundary.

Let $\mathcal C_g$ be the curve complex: vertices are isotopy classes of essential simple closed curves on the closed surface, and simplices are collections admitting mutually disjoint representatives. The mapping class group acts on its geometric realization. Put $Y_g=|\mathcal C_g|/\operatorname{Mod}_g$ and write
\[
 C_{\mathrm{top}}(Y_g)=(Y_g\times[0,\infty))/(Y_g\times\{0\})
\]
for its open, unbounded topological cone, including the apex and with no terminal base face.

\textbf{Conjecture 4.7 ($\Q$-rank of moduli space).}
\[
 \operatorname{Cone}(\mathcal M_g)\ \cong\ C_{\mathrm{top}}(Y_g)
\]
as topological spaces. The source proposes a stronger version specifying the precise bilipschitz type of the natural metric on the cone, but the excerpt supplies neither that metric model on $C_{\mathrm{top}}(Y_g)$ nor quantitative constants. It cites Hattori [Hat] for an analogous statement for locally symmetric lattice quotients $\Gamma\backslash G/K$.

For a genus-$g$ surface $S$ with $n$ punctures, the surrounding local record 11000043 sets $d(S)=3g-3+n$. The source predicts corresponding cone dimensions $\Q$-$\operatorname{rank}(\Gamma)$ and $d(S)$, respectively; in the closed case the latter is $3g-3$. Its analogies pair the root lattice with simple closed curves, simple roots with topological types of simple closed curves, and $\operatorname{Cone}(\Gamma\backslash G/K)$ with $\operatorname{Cone}(\mathcal M_g)$.

The author reports work with H. Masur on a candidate coarse fundamental domain described by length inequalities analogous to root inequalities for Weyl chambers; no inequalities are supplied here. The source also suggests an implication, together with methods in [Ch], to the second statement of Conjecture 4.5, whose target is not reproduced in this record.
\subsection{Short English statement}
The large-scale pointed metric limit of moduli space with the Teichm\"uller metric is conjectured to be the open cone on the mapping-class quotient of the curve complex. The source also proposes a stronger metric description and a dimension analogy with arithmetic $\Q$-rank.
\subsection{Sources}
\begin{itemize}
\item Local UnsolvedMath registry, permanent ID 11000045, source code catalogue entry; no individual proposer or proposal year supplied. Adjacent local IDs 11000044 and 11000043 supply the pointed cone definition, curve complex, and $d(S)$ convention.
\item Farb, edited volume \emph{Problems on Mapping Class Groups and Related Topics}, Conjecture 4.7, PDF page 40; \url{https://www.math.uchicago.edu/~farb/papers/mcgbook.pdf}. The registry records PDF-font-structure extraction and access on 2026-07-29; the inventory has 329 labeled headings distinguished by source page.
\item The excerpt names H. Masur and source references Hattori [Hat] and [Ch]. Full bibliography entries, coarse-domain inequalities, and the referenced second statement of Conjecture 4.5 are not supplied.
\end{itemize}
\end{document}
